
\section{Backlog del Producto}

\begin{xltabular}{\textwidth}{|p{2.2cm}|X|p{1.4cm}|p{1.7cm}|p{1.5cm}|}
\hline
\rowcolor{azulPastelApagado}
\textbf{Ident.} & \textbf{Título} & \textbf{Estim.} & \textbf{Iteración} & \textbf{Entrega} \\ \hline
\endhead

\hline
\endfoot

\rowcolors{2}{azulclaro}{white}
HU-P1 & Como profesor/a, quiero iniciar sesión en la aplicación para consultar la agenda escolar. & 3 & 1 & 1 \\ \hline
HU-P2 & Como profesor/a, quiero consultar las actividades programadas y sus detalles para conocer la información necesaria sobre cada actividad. & 2 & 1 & 1 \\ \hline
HU-P3 & Como profesor/a, quiero crear actividades para mantener actualizada la planificación escolar según las necesidades del alumnado. & 5 & 1 & 1 \\ \hline
HU-P6 & Como profesor/a, quiero asignar una actividad a un estudiante o grupo para adaptar la planificación a sus necesidades educativas. & 3 & 1 & 1 \\ \hline
HU-E1 & Como estudiante, quiero iniciar sesión en la aplicación, para acceder a mi agenda y a mis actividades personales. & 3 & 1 & 1 \\ \hline
HU-E2 & Como estudiante, quiero consultar mi agenda personalizada, para saber qué actividades tengo programadas durante el día. & 3 & 1 & 1 \\ \hline
HU-E3 & Como estudiante, quiero consultar la información de una actividad para comprender fácilmente qué tengo que hacer. & 3 & 1 & 1 \\ \hline
HU-E4 & Como estudiante, quiero consultar las actividades de los próximos días, para organizarme y saber qué tengo programado durante la semana. & 2 & 1 & 1 \\ \hline
HUT-A0 & Como administrador quiero montar infraestructura de la aplicación para que mis usuarios puedan utilizarla. & 2 & 1 & 1 \\ \hline
HU-A1 & Como administrador quiero crear un usuario estudiante para que cierto estudiante pueda usar la aplicación. & 2 & 1 & 1 \\ \hline
HU-A2 & Como administrador quiero crear usuarios profesor y familia para que cierto profesor y familiar pueda usar la aplicación & 1 & 1 & 1 \\ \hline
HU-A4 & Como administrador quiero asignar un estudiante /profesor a un grupo de clase para que yo y los usuarios trabajen con interfaces con menos opciones & 2 & 1 & 1 \\ \hline
HU-A5 & Como administrador quiero asignar un profesor a un estudiante para que el profesor pueda interactuar con la agenda del estudiante & 2 & 1 & 1 \\ \hline
HU-F1 & Como Familia quiero iniciar sesión en la aplicación para acceder a la funcionalidad de esta. & 3 & 2 & 2 \\ \hline
HU-F2 & Como Familia quiero seleccionar uno de mis hijos para acceder a su perfil. & 1 & 2 & 2 \\ \hline
HU-F5 & Como Familia quiero comprobar la agenda de un hijo para ver sus actividads planeadas. & 2 & 2 & 2 \\ \hline
HU-F7 & Como Familia quiero crear una actividad para un hijo para que pueda tenerla presente y pueda realizarla por su cuenta. & 5 & 2 & 2 \\ \hline
HU-P4 & Como profesor/a, quiero cancelar actividades para mantener actualizada la planificación escolar según las necesidades del alumnado. & 5 & 2 & 2\\ \hline
HU-P7 & Como profesor/a, quiero consultar qué estudiantes participan en una actividad para conocer quiénes deben realizarla y organizar su desarrollo. & 2 & 2 & 2 \\ \hline
HU-P10 & Como profesor/a, quiero consultar la información relevante de un estudiante para adaptar las actividades a sus necesidades y facilitar su participación y aprendizaje. & 3 & 2 & 2 \\ \hline
HU-E7 & Como estudiante, quiero marcar los pasos de una actividad conforme los voy realizando, para saber qué pasos he completado y cuál debo realizar a continuación. & 2 & 2 & 2 \\ \hline
HU-E14 & Como estudiante, quiero adaptar las preferencias de visualización de mi agenda a mis necesidades, para consultar las actividades de una forma que me resulte cómoda y comprensible. & 8 & 2 & 2 \\ \hline
HU-A6 & Como administrador quiero asignar un padre/madre/tutor a un estudiante para que el padre pueda/madre/tutor pueda interactuar con la agenda del estudiante. & 2 & 2 & 2 \\ \hline
HU-F3 & Como Familia quiero consultar el perfil de un hijo para ver su configuración. & 1 & 3 & 3 \\ \hline
HU-F4 & Como Familia quiero cambiar la configuración de un hijo para hacer la aplicación más accesible para él/ella. & 1 & 3 & 3 \\ \hline
HU-F6 & Como Familia quiero consultar los detalles específicos de una actividad que tiene mi hijo para conocer los detalles de esta. & 2 & 3 & 3 \\ \hline
HU-F8 & Como Familia quiero editar una actividad que le cree a mi hijo para corregir o añadir detalles de esta. & 5 & 3 & 3 \\ \hline
HU-F9 & Como Familia quiero borrar una actividad que le cree a mi hijo para que no la realice. & 5 & 3 & 3 \\ \hline
HU-P5 & Como profesor/a, quiero editar actividades para mantener actualizada la planificación escolar según las necesidades del alumnado. & 5 & 3 & 3 \\ \hline
HU-P9 & Como profesor/a, quiero marcar una actividad como realizada para mantener actualizado el calendario y registrar las actividades que ya se han llevado a cabo. & 2 & 3 & 3 \\ \hline
HU-E5 & Como estudiante, quiero añadir actividades personales a mi agenda, para organizar mis actividads y actividades cotidianas en casa de forma autónoma. & 5 & 3 & 3 \\ \hline
HU-E8 & Como estudiante, quiero consultar el tiempo restante de una actividad, para saber cuánto tiempo tengo para realizarla. & 2 & 3 & 3 \\ \hline
HU-E9 & Como estudiante, quiero confirmar el inicio de una actividad, para registrar su inicio y llevar un seguimiento de mi progreso. & 1 & 3 & 3 \\ \hline
HU-E10 & Como estudiante, quiero confirmar la finalización de una actividad, para registrar su finalización y saber qué actividad debo realizar a continuación. & 1 & 3 & 3 \\ \hline
HU-A3 & Como administrador quiero eliminar usuarios para cuando dejen de estar en la institución y ahorrar almacenamiento. & 1 & 3 & 3 \\ \hline
HU-P8 & Como profesor/a, quiero consultar qué estudiantes han comenzado o completado una actividad para realizar un seguimiento individual de su progreso. & 8 & 4 & 4 \\ \hline
HU-E6 & Como estudiante, quiero recibir avisos cuando llegue la hora de realizar una actividad, para saber cuándo debo comenzar. & 8 & 4 & 4 \\ \hline
HU-E11 & Como estudiante, quiero recibir un aviso cuando se cancele una actividad de mi agenda, para saber que ya no tengo que realizarla. & 3 & 4 & 4 \\ \hline
HU-F10 & Como Familia quiero abrir un canal de comunicación en la aplicación con mi hijo para poder comunicarme con él & 21 & 5 & 5 \\ \hline
HU-E12 & Como estudiante, quiero recibir una sugerencia de actividad alternativa cuando se cancele una actividad programada, para saber qué puedo hacer en su lugar y mantener organizada mi rutina. & 21 & 5 & 5 \\ \hline
HU-E13 & Como estudiante, quiero valorar una actividad realizada, para que la aplicación pueda tener en cuenta mis preferencias en futuras sugerencias. & 5 & 5 & 5 \\ \hline
HU-E15 & Como estudiante, quiero enviar mensajes a través de un chat accesible, para solicitar ayuda o resolver dudas mientras realizo mis actividades. & 21 & 5 & 5 \\ \hline
HU-E16 & Como estudiante, quiero recibir mensajes a través de un chat accesible, para recibir indicaciones, en caso de necesitarlo, mientras realizo mis actividades. & 21 & 5 & 5 \\ \hline
\end{xltabular}
\noindent\textit{Nota sobre el alcance de las iteraciones:} El desarrollo principal del
 proyecto abarca estrictamente las **Iteraciones 1, 2 y 3**. Las Historias de Usuario asignadas 
 a la **Iteración 4** únicamente se abordarán si se dispone de tiempo suficiente tras completar las iteraciones principales en la entrega final. Por su parte, las funcionalidades marcadas en la **Iteración 5** quedan explícitamente fuera del alcance comprometido para este proyecto debido a su elevada complejidad.
\vspace{0.4cm}
\subsection{Tarjetas de la HU (Primera Iteración)}

\begin{table}[H]
\centering
\rowcolors{1}{white}{white}
\begin{tabularx}{\textwidth}{|X|X|X|}
\hline
\rowcolor{azulclaro}
\multicolumn{1}{|l|}{\textbf{Identificador:} HU-P1} & \multicolumn{2}{l|}{\textbf{Iniciar sesión}} \\ \hline
\multicolumn{3}{|p{\dimexpr\textwidth-2\tabcolsep-2\arrayrulewidth}|}{\textbf{Descripción:} Como profesor/a, quiero iniciar sesión en la aplicación para consultar la agenda escolar.} \\ \hline
\rowcolor{white}
\textbf{Estimación: 3} & \textbf{Prioridad: 1} & \textbf{Entrega: 1} \\ \hline
\multicolumn{3}{|p{\dimexpr\textwidth-2\tabcolsep-2\arrayrulewidth}|}{%
\textbf{Pruebas de aceptación:}
\begin{itemize}[leftmargin=1.5em, itemsep=2pt, topsep=3pt]
    \item En caso de introducir unas credenciales válidas, el sistema permite al profesor/a iniciar sesión y acceder a la aplicación.
    \item En caso de introducir un correo electrónico o una contraseña incorrectos, el sistema avisa de que las credenciales no son válidas y deniega el acceso.
    \item En caso de intentar iniciar sesión sin introducir alguno de los campos obligatorios, el sistema avisa de que faltan datos y deniega el acceso.
    \item En caso de intentar iniciar sesión con una cuenta que no está registrada en el sistema, el sistema avisa de que las credenciales no son válidas y deniega el acceso.
\end{itemize}
} \\ \hline
\multicolumn{3}{|p{\dimexpr\textwidth-2\tabcolsep-2\arrayrulewidth}|}{\textbf{Observaciones:} Requiere implementar un sistema de autenticación de usuarios.} \\ \hline
\end{tabularx}
\end{table}

\vspace{0.4cm}

\begin{table}[H]
\centering
\rowcolors{1}{white}{white}
\begin{tabularx}{\textwidth}{|X|X|X|}
\hline
\rowcolor{azulclaro}
\multicolumn{1}{|l|}{\textbf{Identificador:} HU-P2} & \multicolumn{2}{l|}{\textbf{Consultar actividades programadas y sus detalles}} \\ \hline
\multicolumn{3}{|p{\dimexpr\textwidth-2\tabcolsep-2\arrayrulewidth}|}{\textbf{Descripción:} Como profesor/a, quiero consultar las actividades programadas y sus detalles para conocer la información necesaria sobre cada actividad.} \\ \hline
\rowcolor{white}
\textbf{Estimación: 2} & \textbf{Prioridad: 1} & \textbf{Entrega: 1} \\ \hline
\multicolumn{3}{|p{\dimexpr\textwidth-2\tabcolsep-2\arrayrulewidth}|}{%
\textbf{Pruebas de aceptación:}
\begin{itemize}[leftmargin=1.5em, itemsep=2pt, topsep=3pt]
    \item En caso de seleccionar una actividad existente, el sistema muestra sus detalles, incluyendo la fecha, la hora, la descripción y el alumnado al que está dirigida, cuando corresponda.
    \item En caso de consultar un día en el que no hay actividades programadas, el sistema informa de que no hay actividades previstas.
    \item En caso de intentar consultar una actividad inexistente o que ya no está disponible, el sistema avisa de que no se ha encontrado la actividad.
    \item En caso de que una actividad haya sido modificada, el sistema muestra la información actualizada de dicha actividad.
\end{itemize}
} \\ \hline
\multicolumn{3}{|p{\dimexpr\textwidth-2\tabcolsep-2\arrayrulewidth}|}{\textbf{Observaciones:} Requiere implementar el sistema de actividades y su consulta desde la agenda escolar.} \\ \hline
\end{tabularx}
\end{table}

\vspace{0.4cm}

\begin{table}[H]
\centering
\rowcolors{1}{white}{white}
\begin{tabularx}{\textwidth}{|X|X|X|}
\hline
\rowcolor{azulclaro}
\multicolumn{1}{|l|}{\textbf{Identificador:} HU-P3} & \multicolumn{2}{l|}{\textbf{Crear actividades}} \\ \hline
\multicolumn{3}{|p{\dimexpr\textwidth-2\tabcolsep-2\arrayrulewidth}|}{\textbf{Descripción:} Como profesor/a, quiero crear actividades para mantener actualizada la planificación escolar según las necesidades del alumnado} \\ \hline
\rowcolor{white}
\textbf{Estimación: 5} & \textbf{Prioridad: 1} & \textbf{Entrega: 1} \\ \hline
\multicolumn{3}{|p{\dimexpr\textwidth-2\tabcolsep-2\arrayrulewidth}|}{%
\textbf{Pruebas de aceptación:}
\begin{itemize}[leftmargin=1.5em, itemsep=2pt, topsep=3pt]
    \item En caso de intentar crear una actividad con la información necesaria, el sistema registra la actividad para su posterior asignación.
    \item En caso de intentar crear una actividad pero falta algún dato necesario, el sistema avisa de la falta de datos obligatorios y previene registrar la actividad.
    \item En caso de intentar crear una tarea con algún dato incorrecto, el sistema avisa de los datos incorrectos y previene registrar la actividad.
\end{itemize}
} \\ \hline
\multicolumn{3}{|p{\dimexpr\textwidth-2\tabcolsep-2\arrayrulewidth}|}{\textbf{Observaciones:} Requiere crear el sistema de actividades~} \\ \hline
\end{tabularx}
\end{table}

\vspace{0.4cm}

\begin{table}[H]
\centering
\rowcolors{1}{white}{white}
\begin{tabularx}{\textwidth}{|X|X|X|}
\hline
\rowcolor{azulclaro}
\multicolumn{1}{|l|}{\textbf{Identificador:} HU-P6} & \multicolumn{2}{l|}{\textbf{Asignar actividad}} \\ \hline
\multicolumn{3}{|p{\dimexpr\textwidth-2\tabcolsep-2\arrayrulewidth}|}{\textbf{Descripción:} Como profesor/a, quiero asignar una actividad a un estudiante o grupo para adaptar la planificación a sus necesidades educativas.~} \\ \hline
\rowcolor{white}
\textbf{Estimación: 3} & \textbf{Prioridad: 1} & \textbf{Entrega: 1} \\ \hline
\multicolumn{3}{|p{\dimexpr\textwidth-2\tabcolsep-2\arrayrulewidth}|}{%
\textbf{Pruebas de aceptación:}
\begin{itemize}[leftmargin=1.5em, itemsep=2pt, topsep=3pt]
    \item En caso de asignar la actividad a un estudiante o selección de estos registrados, el sistema realiza la asignación.
    \item En caso de asignar la actividad a un estudiante o selección de estos pero alguno no existe, el sistema avisa y rechaza la asignación.
    \item En caso de asignar la actividad a un grupo de estudiantes existente, el sistema realiza la asignación.
    \item En caso de asignar la actividad a un grupo de estudiantes no existente, el sistema avisa y rechaza la asignación.
    \item En caso de intentar asignar una actividad inexistente, el sistema avisa y rechaza la asignación.
\end{itemize}
} \\ \hline
\multicolumn{3}{|p{\dimexpr\textwidth-2\tabcolsep-2\arrayrulewidth}|}{\textbf{Observaciones:} Requiere crear el sistema de usuarios.} \\ \hline
\end{tabularx}
\end{table}

\vspace{0.4cm}

\begin{table}[H]
\centering
\rowcolors{1}{white}{white}
\begin{tabularx}{\textwidth}{|X|X|X|}
\hline
\rowcolor{azulclaro}
\multicolumn{1}{|l|}{\textbf{Identificador:} HU-E1} & \multicolumn{2}{l|}{\textbf{Iniciar sesión estudiante}} \\ \hline
\multicolumn{3}{|p{\dimexpr\textwidth-2\tabcolsep-2\arrayrulewidth}|}{\textbf{Descripción:} Como estudiante, quiero iniciar sesión en la aplicación, para acceder a mi agenda y a mis actividades personales.} \\ \hline
\rowcolor{white}
\textbf{Estimación: 3} & \textbf{Prioridad: 1} & \textbf{Entrega: 1} \\ \hline
\multicolumn{3}{|p{\dimexpr\textwidth-2\tabcolsep-2\arrayrulewidth}|}{%
\textbf{Pruebas de aceptación:}
\begin{itemize}[leftmargin=1.5em, itemsep=2pt, topsep=3pt]
    \item En caso de introducir unas credenciales válidas, el sistema permite al estudiante /a iniciar sesión y acceder a la aplicación.
    \item En caso de introducir un correo electrónico o una contraseña incorrectos, el sistema avisa de que las credenciales no son válidas y deniega el acceso.
    \item En caso de intentar iniciar sesión sin introducir alguno de los campos obligatorios, el sistema avisa de que faltan datos y deniega el acceso.
    \item En caso de intentar iniciar sesión con una cuenta que no está registrada en el sistema, el sistema avisa de que las credenciales no son válidas y deniega el acceso.
\end{itemize}
} \\ \hline
\multicolumn{3}{|p{\dimexpr\textwidth-2\tabcolsep-2\arrayrulewidth}|}{\textbf{Observaciones:} Requiere crear el sistema de usuarios y actividades.} \\ \hline
\end{tabularx}
\end{table}

\vspace{0.4cm}

\begin{table}[H]
\centering
\rowcolors{1}{white}{white}
\begin{tabularx}{\textwidth}{|X|X|X|}
\hline
\rowcolor{azulclaro}
\multicolumn{1}{|l|}{\textbf{Identificador:} HU-E2} & \multicolumn{2}{l|}{\textbf{Consultar agenda}} \\ \hline
\multicolumn{3}{|p{\dimexpr\textwidth-2\tabcolsep-2\arrayrulewidth}|}{\textbf{Descripción:} Como estudiante, quiero consultar mi agenda personalizada, para saber qué actividades tengo programadas durante el día.} \\ \hline
\rowcolor{white}
\textbf{Estimación: 3} & \textbf{Prioridad: 1} & \textbf{Entrega: 1} \\ \hline
\multicolumn{3}{|p{\dimexpr\textwidth-2\tabcolsep-2\arrayrulewidth}|}{%
\textbf{Pruebas de aceptación:}
\begin{itemize}[leftmargin=1.5em, itemsep=2pt, topsep=3pt]
    \item La interfaz de la agenda se muestra tal como indican sus preferencias personales.
    \item En caso de tener actividades en la agenda, que aparezcan listadas por pantalla.
    \item En caso de que su perfil lo permita, poder cambiar entre la vista semanal y la diaria.
\end{itemize}
} \\ \hline
\multicolumn{3}{|p{\dimexpr\textwidth-2\tabcolsep-2\arrayrulewidth}|}{\textbf{Observaciones:}} \\ \hline
\end{tabularx}
\end{table}

\vspace{0.4cm}

\begin{table}[H]
\centering
\rowcolors{1}{white}{white}
\begin{tabularx}{\textwidth}{|X|X|X|}
\hline
\rowcolor{azulclaro}
\multicolumn{1}{|l|}{\textbf{Identificador:} HU-E3} & \multicolumn{2}{l|}{\textbf{Consultar la información de una actividad}} \\ \hline
\multicolumn{3}{|p{\dimexpr\textwidth-2\tabcolsep-2\arrayrulewidth}|}{\textbf{Descripción:} Como estudiante, quiero consultar la información de una actividad para comprender fácilmente qué tengo que hacer} \\ \hline
\rowcolor{white}
\textbf{Estimación: 3} & \textbf{Prioridad: 1} & \textbf{Entrega: 1} \\ \hline
\multicolumn{3}{|p{\dimexpr\textwidth-2\tabcolsep-2\arrayrulewidth}|}{%
\textbf{Pruebas de aceptación:}
\begin{itemize}[leftmargin=1.5em, itemsep=2pt, topsep=3pt]
    \item En caso que la actividad esté dentro del horario programado y pendiente de realizar, cuando el estudiante la selecciona, el sistema muestra la vista detallada con la información y las instrucciones para poder llevarla a cabo.
    \item En caso que la actividad ya haya pasado de su hora programada o se encuentre vencida, cuando el estudiante intenta consultarla, el sistema muestra un aviso indicando que el tiempo ha expirado y muestra los detalles en modo lectura.
    \item En caso que la actividad ya se haya dado por finalizada o completada, cuando el estudiante intenta consultarla, el sistema despliega la información con una marca visual de completada.
\end{itemize}
} \\ \hline
\multicolumn{3}{|p{\dimexpr\textwidth-2\tabcolsep-2\arrayrulewidth}|}{\textbf{Observaciones:} Será relativo a que tareas expiran y cuando.} \\ \hline
\end{tabularx}
\end{table}

\vspace{0.4cm}

\begin{table}[H]
\centering
\rowcolors{1}{white}{white}
\begin{tabularx}{\textwidth}{|X|X|X|}
\hline
\rowcolor{azulclaro}
\multicolumn{1}{|l|}{\textbf{Identificador:} HU-E4} & \multicolumn{2}{l|}{\textbf{Consultar actividades}} \\ \hline
\multicolumn{3}{|p{\dimexpr\textwidth-2\tabcolsep-2\arrayrulewidth}|}{\textbf{Descripción:} Como estudiante, quiero consultar las actividades de los próximos días, para organizarme y saber qué tengo programado durante la semana.} \\ \hline
\rowcolor{white}
\textbf{Estimación: 3} & \textbf{Prioridad: 1} & \textbf{Entrega: 1} \\ \hline
\multicolumn{3}{|p{\dimexpr\textwidth-2\tabcolsep-2\arrayrulewidth}|}{%
\textbf{Pruebas de aceptación:}
\begin{itemize}[leftmargin=1.5em, itemsep=2pt, topsep=3pt]
    \item En caso que la agenda cuente con más actividades, el sistema muestra una vista semanal con las tareas por venir en formato calendario semanal.
    \item En caso que la agenda no tenga actividades, el sistema muestra un mensaje avisando.
\end{itemize}
} \\ \hline
\multicolumn{3}{|p{\dimexpr\textwidth-2\tabcolsep-2\arrayrulewidth}|}{\textbf{Observaciones:}} \\ \hline
\end{tabularx}
\end{table}

\vspace{0.4cm}

\begin{table}[H]
\centering
\rowcolors{1}{white}{white}
\begin{tabularx}{\textwidth}{|X|X|X|}
\hline
\rowcolor{azulclaro}
\multicolumn{1}{|l|}{\textbf{Identificador:} HUT-A0} & \multicolumn{2}{l|}{\textbf{Montar infraestructura de la web}} \\ \hline
\multicolumn{3}{|p{\dimexpr\textwidth-2\tabcolsep-2\arrayrulewidth}|}{\textbf{Descripción:} Como administrador quiero montar infraestructura de la web para que mis usuarios puedan utilizarla.} \\ \hline
\rowcolor{white}
\textbf{Estimación: 2} & \textbf{Prioridad: 1} & \textbf{Entrega: 1} \\ \hline
\multicolumn{3}{|p{\dimexpr\textwidth-2\tabcolsep-2\arrayrulewidth}|}{%
\textbf{Pruebas de aceptación:}
\begin{itemize}[leftmargin=1.5em, itemsep=2pt, topsep=3pt]
    \item En caso que la base de datos o los servicios requeridos no estén iniciados, cuando se intenta arrancar la infraestructura de la aplicación, el sistema detiene el despliegue y notifica la falta de dependencias activas.
    \item En caso que los puertos o recursos del servidor estén ocupados por otro proceso, cuando se intenta levantar la infraestructura, el sistema genera un registro de error.
    \item En caso que todos los servicios y variables de entorno estén correctamente configurados, cuando se realiza la petición de arranque de la infraestructura, el sistema despliega el entorno de desarrollo y habilita la conexión con la base de datos.
\end{itemize}
} \\ \hline
\multicolumn{3}{|p{\dimexpr\textwidth-2\tabcolsep-2\arrayrulewidth}|}{\textbf{Observaciones:} Es una tarea técnica y la base de todo para la ejecución de pruebas.} \\ \hline
\end{tabularx}
\end{table}

\vspace{0.4cm}

\begin{table}[H]
\centering
\rowcolors{1}{white}{white}
\begin{tabularx}{\textwidth}{|X|X|X|}
\hline
\rowcolor{azulclaro}
\multicolumn{1}{|l|}{\textbf{Identificador:} HU-A1} & \multicolumn{2}{l|}{\textbf{Crear usuario estudiante}} \\ \hline
\multicolumn{3}{|p{\dimexpr\textwidth-2\tabcolsep-2\arrayrulewidth}|}{\textbf{Descripción:} Como administrador quiero crear un usuario estudiante para que cierto estudiante pueda usar la aplicación} \\ \hline
\rowcolor{white}
\textbf{Estimación: 2} & \textbf{Prioridad: 1} & \textbf{Entrega: 1} \\ \hline
\multicolumn{3}{|p{\dimexpr\textwidth-2\tabcolsep-2\arrayrulewidth}|}{%
\textbf{Pruebas de aceptación:}
\begin{itemize}[leftmargin=1.5em, itemsep=2pt, topsep=3pt]
    \item En caso de introducir todos los datos obligatorios con información válida, el sistema crea el usuario estudiante y confirma que se ha registrado correctamente.
    \item En caso de intentar crear un usuario estudiante sin introducir alguno de los campos obligatorios, el sistema avisa de que faltan datos y no crea el usuario.
    \item En caso de introducir un usuario que ya está registrado,el sistema avisa de que ya existe ese usuario y no crea la cuenta duplicada.
    \item En caso de introducir datos con un formato no válido en alguno de los campos obligatorios, el sistema avisa del error y no crea al usuario.
\end{itemize}
} \\ \hline
\multicolumn{3}{|p{\dimexpr\textwidth-2\tabcolsep-2\arrayrulewidth}|}{\textbf{Observaciones:} Requiere disponer de un sistema de gestión de usuarios y de los datos necesarios para registrar a un alumno.} \\ \hline
\end{tabularx}
\end{table}

\vspace{0.4cm}

\begin{table}[H]
\centering
\rowcolors{1}{white}{white}
\begin{tabularx}{\textwidth}{|X|X|X|}
\hline
\rowcolor{azulclaro}
\multicolumn{1}{|l|}{\textbf{Identificador:} HU-A2} & \multicolumn{2}{l|}{\textbf{Crear profesores y familiares.}} \\ \hline
\multicolumn{3}{|p{\dimexpr\textwidth-2\tabcolsep-2\arrayrulewidth}|}{\textbf{Descripción:} Como administrador quiero crear un usuario profesor/tutor para que cierto profesor/tutor pueda usar la aplicación} \\ \hline
\rowcolor{white}
\textbf{Estimación: 1} & \textbf{Prioridad: 1} & \textbf{Entrega: 1} \\ \hline
\multicolumn{3}{|p{\dimexpr\textwidth-2\tabcolsep-2\arrayrulewidth}|}{%
\textbf{Pruebas de aceptación:}
\begin{itemize}[leftmargin=1.5em, itemsep=2pt, topsep=3pt]
    \item En caso de que no exista un usuario con el correo de creación, se crea dicho usuario en el sistema.
    \item En caso de que exista el usuario que se quiere crear, avisar al administrador de su existencia y no dejar que se cree un duplicado.
    \item El criterio de discriminación entre profesor y familia funciona correctamente, un profesor es un profesor y un familiar es un familiar.
\end{itemize}
} \\ \hline
\multicolumn{3}{|p{\dimexpr\textwidth-2\tabcolsep-2\arrayrulewidth}|}{\textbf{Observaciones:}} \\ \hline
\end{tabularx}
\end{table}

\vspace{0.4cm}

\begin{table}[H]
\centering
\rowcolors{1}{white}{white}
\begin{tabularx}{\textwidth}{|X|X|X|}
\hline
\rowcolor{azulclaro}
\multicolumn{1}{|l|}{\textbf{Identificador:} HU-A4} & \multicolumn{2}{l|}{\textbf{Asignar estudiante y profesor a grupos}} \\ \hline
\multicolumn{3}{|p{\dimexpr\textwidth-2\tabcolsep-2\arrayrulewidth}|}{\textbf{Descripción:} Como administrador quiero asignar un estudiante /profesor a un grupo de clase para que yo y los usuarios trabajen con interfaces con menos opciones} \\ \hline
\rowcolor{white}
\textbf{Estimación: 1} & \textbf{Prioridad: 1} & \textbf{Entrega: 1} \\ \hline
\multicolumn{3}{|p{\dimexpr\textwidth-2\tabcolsep-2\arrayrulewidth}|}{%
\textbf{Pruebas de aceptación:}
\begin{itemize}[leftmargin=1.5em, itemsep=2pt, topsep=3pt]
    \item En caso de que exista el usuario y el grupo, el sistema registra la asociación.
    \item En caso de que no exista el usuario o el grupo, el sistema avisa y rechaza la operación.
    \item En caso de que el usuario ya estuviera en el grupo, el sistema avisa y rechaza la asignación.
\end{itemize}
} \\ \hline
\multicolumn{3}{|p{\dimexpr\textwidth-2\tabcolsep-2\arrayrulewidth}|}{\textbf{Observaciones:}} \\ \hline
\end{tabularx}
\end{table}

\vspace{0.4cm}

\begin{table}[H]
\centering
\rowcolors{1}{white}{white}
\begin{tabularx}{\textwidth}{|X|X|X|}
\hline
\rowcolor{azulclaro}
\multicolumn{1}{|l|}{\textbf{Identificador:} HU-A5} & \multicolumn{2}{l|}{\textbf{Asignar un profesor a un alumno}} \\ \hline
\multicolumn{3}{|p{\dimexpr\textwidth-2\tabcolsep-2\arrayrulewidth}|}{\textbf{Descripción:} Como administrador quiero asignar un profesor a un estudiante para que el profesor pueda interactuar con la agenda del estudiante.} \\ \hline
\rowcolor{white}
\textbf{Estimación: 2} & \textbf{Prioridad: 1} & \textbf{Entrega: 1} \\ \hline
\multicolumn{3}{|p{\dimexpr\textwidth-2\tabcolsep-2\arrayrulewidth}|}{%
\textbf{Pruebas de aceptación:}
\begin{itemize}[leftmargin=1.5em, itemsep=2pt, topsep=3pt]
    \item En caso de seleccionar un profesor y un alumno registrados en el sistema, el sistema realiza la asignación correctamente y confirma que se ha completado.
    \item En caso de intentar realizar una asignación sin seleccionar un profesor o un alumno, el sistema avisa de que faltan datos y no realiza la asignación.
    \item En caso de seleccionar un profesor que no está registrado en el sistema, el sistema avisa de que no se puede realizar la asignación.
    \item En caso de seleccionar un alumno que no está registrado en el sistema, el sistema avisa de que no se puede realizar la asignación.
    \item En caso de asignar un profesor a un alumno correctamente, el profesor puede consultar e interactuar con la agenda de ese alumno según los permisos establecidos.
    \item En caso de intentar asignar de nuevo al mismo profesor un alumno que ya tiene asignado, el sistema evita duplicar la asignación e informa de que ya existe.
\end{itemize}
} \\ \hline
\multicolumn{3}{|p{\dimexpr\textwidth-2\tabcolsep-2\arrayrulewidth}|}{\textbf{Observaciones:} Requiere disponer de usuarios de tipo profesor y alumno previamente registrados, así como de un sistema de gestión de las asignaciones entre ambos.~} \\ \hline
\end{tabularx}
\end{table}